\section{Cost}
\label{sec:complexity}

The current \SPEC{} (3.1.0-rc.2) makes \spec{13.11} normative for semantic
equivalence and informative for cost and technique. An implementation must
produce the plan \spec{13} defines, but need not use or attain any bound
below. The accounting assumes decoded, indexed inputs with fixed-width or
interned identifiers and adjacency and precedence inputs already in their
specified order; comparison, decoding, sorting, allocation, adapter, record
matching and emitted output costs remain, and hash-table operations are
expected constant time under the chosen hashing assumptions. This section
transcribes the phase table, proves the transformations it relies on and the
parser bound \spec{18.3} depends on, and states the space ledger; it makes no
global optimality claim.

Notation follows \spec{13.11} so that the table can be read beside it. $P$
packages, $E$ workspace dependency edges, $H$ commits and $A$ parent edges in
the history reachable from \Head{}, $C$ commits in the union of all pending
windows, $U$ units in those commits, $N$ total bytes of their messages and
changed paths, $T$ reachable tags, and $M$ tag-record and package-format
matches examined while partitioning the inventory (at worst $P \cdot T$). $k$
is the number of distinct commits carrying a boundary: a stable baseline, a
package's newest baseline of any channel where that is another commit, or the
release of a provider that bounds a consumer's owed window (\spec{13.3}), of
which there are none while no consumer has got ahead of a provider. $m$ is the
number of distinct marker commits, the commits some ancestry question of
\spec{13} is asked about: the $k$ boundaries, every \code{cancel} commit, and
every commit carrying an \code{Edits}, \code{Deletes} or \code{Reverts}
footer. $R$ is the work of resolving scopes and changed paths, $I$ the number
of unit-to-package incidences (up to $P \cdot U$), $Z$ the retained or
emitted contribution and provenance incidences and $Z_v$ those examined while
constructing them, $I_w$ the stored or emitted package-window memberships (up
to $P \cdot C$). $F_c$ counts correction footer selectors, $J_{corr}$ the
package-target incidences examined for correction containment and
application, $I_c$ cancel-scope incidences, $P_c$ packages named by a cancel,
and $J_c$ cancellation queries across direct admission and both propagation
axes. $\mathrm{bw}(x) = \max(1, \lceil x / w \rceil)$ for word size $w$, so a
zero-marker phase still pays for traversal. $F$ is the number of publish
failures and $O_{out}$ the size of emitted output. In the per-target row only,
$D$ is the number of targets one unit reaches, $S$ its source-set size and
$\Sigma$ its resolved scope-set size. For the polyrepository profile, $Q$ is
the number of repositories, $H_q$ and $A_q$ one repository's reachable commits
and parent edges, $G$ the control gitlink transitions examined, $K_q$ the
distinct boundary revisions used in repository $q$, $m_q$ the markers in $q$,
$V$ the package memberships in shared-version groups, $X$ the (consumer
release tag, repository) boundaries resolved under \spec{27.11}, $Y$ the
longest route between two peers and $B$ the distinct (repository, revision)
trees read. For the execution profile, $T_k$ is the run's tasks, $D_k$ their
precedence edges, $W_n$ the worker nodes, $P_o$ the packages whose outputs
another task consumes, and $L_r$ the coordination refs alive in one transport
repository.

\begin{center}
\footnotesize
\begin{longtable}{@{}L{3.4cm}L{2.7cm}L{3.2cm}L{4.9cm}@{}}
\toprule
Phase & Literal & Achievable & Note \\
\midrule
\endhead
Load workspace (\spec{13.1}) & $O(P + E)$ & $O(P + E)$ & \\
Load tags (\spec{13.2}) & $O(T \log T + M)$ & $O(T + M)$ & running maximum per package; hash (package, version) for \diag{E191} \\
Pending windows (\spec{13.3}) & $O(P \cdot (H + A))$ & $O((H + A)\,\mathrm{bw}(m) + I_w)$ & one marker pass, \cref{prop:marker-pass}; $O((k+1)(H+A) + I_w)$ with a walk per boundary, \cref{prop:windows} \\
Parse and resolve (\spec{13.4}) & $O(N + R + I)$ & $O(N + R + I)$ & lexing and resolution have different cache keys; \cref{prop:parse} \\
Corrections (\spec{13.4b}) & repeated ancestry and scope scans & $O(H + F_c + J_{corr} + \mathrm{wildcards} \cdot I + O_{out})$ & indexed ancestry charged in the marker pass; package-target probes counted even without wildcards \\
Cancellation predicates (\spec{13.5}, \spec{9.2}) & a cancel scan per query & $O(I_c + (P_c + J_c)\,\mathrm{bw}(\mathrm{cancels}))$ & per-package cancel masks; ancestry charged in the marker pass \\
Direct bumps (\spec{13.6}) & $O(U + I)$ & $O(U + I)$ & consume resolved incidences \\
Holds (\spec{13.6a}) & $O(U + I)$ & $O(U + I)$ & consume resolved incidences \\
Propagation graph walks (\spec{13.7}) & $O(U \cdot (P + E))$ & input-dependent & reuse a walk only when its full inputs match; \cref{thm:g1}, \cref{prop:bucket} \\
Propagation materialisation & repeated set unions & $O(Z_v + Z)$ & in-place indexed insertion; only unique required output has the lower bound $\Omega(Z)$ \\
Per-target predicates (\spec{9.3a}) & $O(D \cdot (S + \Sigma))$ & $O(D + S + \Sigma)$ & per unit; hoist the source-channel set and scope; \cref{prop:hoist} \\
Channel resolution (\spec{13.8}) & $O(P \cdot U)$ & $O(U + I + P)$ & invert unit-to-package incidences; \cref{prop:invert} \\
Versions and plan (\spec{13.9}, \spec{13.10}) & repeated member and record scans & $O(P + I + Z + O_{out})$ & consume built aggregates; scan each disjoint version group once \\
Publish order (\spec{19.2}) & $O(P^2 + E)$ & $O(E + P \log P)$ & a comparison heap over the ready set; \cref{prop:heap} \\
Blocking closure (\spec{19.3}) & $O(P \cdot (P + E))$ & $O(P + E)$ per run & one multi-source reverse traversal; \cref{prop:closure} \\
Build readiness (\spec{19.2a}) & $O(P \cdot (P + E))$ & $O(P + E)$ & a pass-through node per non-building package, or a memoised index; \cref{prop:readiness-cost} \\
Polyrepository snapshots (\spec{27}) & repeated control scans & $O(G + \sum_q (H_q + A_q))$ & index control gitlinks once; walk each source snapshot once \\
Polyrepository windows (\spec{27}) & $O(P \cdot \sum_q (H_q + A_q))$ & $O(\sum_q (H_q + A_q)\,\mathrm{bw}(m_q) + I_w)$ & one marker pass per repository; $O(\sum_q K_q (H_q + A_q) + I_w)$ with a walk per boundary \\
Publication input closure (\spec{27.2}) & $O(P \cdot (P + E + V))$ & $O((P + E + V)\lceil Q / w\rceil)$ & condense the group-augmented graph, one bitset union per edge \\
Link evidence (\spec{27.11}) & $O(X \cdot (Q + Y))$, $X \cdot Y$ tree reads & $O(Q + X \cdot Y)$, $B$ tree reads & root the link tree once; read each (repository, revision) tree once \\
Link settlement (\spec{27.11}) & one commit per route hop & one commit per recording repository & at most $Q - 1$ commits, two in a star \\
Record comparison (\spec{13.2}) & $O(T)$ per repository & $O(T)$ per repository & one inventory read under the lock; \cref{def:records-lock} \\
Task graph and readiness (\spec{28.3}) & $O(T_k \cdot (T_k + D_k))$ & $O(T_k + D_k)$ & indegree counters decrement each successor once \\
Placement (\spec{28.2}) & $O(T_k \cdot W_n)$ & $O(T_k \cdot W_n)$ & $W_n$ is small; a free list per compatibility class is rarely worth its bookkeeping \\
Output transfer (\spec{28.5}) & up to $T_k \cdot P_o$ transfers & at most $P_o \cdot W_n$ transfers & one per distinct (output set, node); verified bytes are not fetched again \\
Coordination polling (\spec{28.4}) & $O(L_r)$ refs per tick & refs addressed to the polling node & the name prefix routes; one fetch per changed tip; dead runs tax every tick until removed \\
\bottomrule
\end{longtable}
\end{center}

Multiplicative costs, corrections, cancellation, provenance and retained
indexes can each dominate and must be charged even when their final result is
empty. The candidate-traversal row is not a bound on the complete planner: a
cached walk is reusable only when its source set, depth and edge kinds are
identical, and it must retain each unit's own distance, fresh-window
membership, cancellation, channel eligibility, self-source exclusion and
provenance. Per-unit predicate checks and output construction add work, and
this report does not prove that all such work fits the bucket traversal
bound. The polyrepository and execution rows are sums over repositories and
counts of tasks respectively; neither is a globally linear end-to-end claim.

\subsection{Windows}

\begin{proposition}[window bucketing]
\label{prop:windows}
$\{\Win{\st}{p} : p \in \Pkgs\}$ has at most $k_s + 1$ distinct members, where
$k_s \le k$ is the number of distinct defined values of $\stcommit{\st}{\cdot}$.
Computing $\Rch{\cdot}$ once per distinct boundary commit, once more for
packages with no stable baseline, and testing membership by lookup costs
$O((k+1)(H + A))$ reachability work and $O(I_w)$ for the memberships stored
or emitted.
\end{proposition}

\begin{proof}
By \cref{def:window}, $\Win{\st}{p}$ depends on $p$ only through
$\stcommit{\st}{p}$; packages sharing a stable baseline commit share a window,
and packages with no stable baseline share $\Rch{\Head}$. So the map from the
$k_s$ stable boundary commits, plus the undefined case, onto distinct windows
is surjective, giving the count. Each traversal visits at most $H$ commits and
$A$ parent edges, there are at most $k + 1$ of them when every boundary of
\spec{13.3}, newest-baseline and owed-window boundaries included, gets one,
and each membership stored or emitted is charged once in $I_w$.
\end{proof}

The traversals range over the reachable history, so the bound is in $H + A$
and not in $C$: proving that a commit is outside a window may require walking
commits that never enter the union. Commit-graph generation numbers can
reject candidates and bound a walk but cannot prove ancestry across branches;
a positive answer still needs a reachability query or a precomputed index. $k$
counts distinct boundary commits, not release runs, and can exceed $P$, since
a release may record packages at different commits and a provider's releases
can bound several consumers' owed windows.

\begin{proposition}[one marker pass, \spec{13.11}]
\label{prop:marker-pass}
Give each of the $m$ markers one bit, and visit the reachable history once,
children before parents: when a commit is visited, set its own marker bit if
it is a marker, then OR its mask into the mask of each of its parents. When
the pass ends, bit $x$ of $\mathrm{mask}(c)$ is set exactly when $c \in
\Rch{x}$. The pass costs $O((H + A)\,\mathrm{bw}(m))$ word operations, against
$O(m \cdot (H + A))$ for a walk per marker, and dense rows occupy $H \cdot
\lceil m / w \rceil$ words.
\end{proposition}

\begin{proof}
$c \in \Rch{x}$ exactly when $c = x$ or some child $c'$ of $c$ satisfies $c'
\in \Rch{x}$, since a path from $x$ down to $c \ne x$ ends with an edge from a
child of $c$ to $c$. Claim: when a commit is visited its mask is final. By
induction along the visiting order: every child of $c$ is visited before $c$,
each child's mask is final at its visit by the induction hypothesis, and each
child ORs that final mask into $\mathrm{mask}(c)$ at its visit; nothing else
writes $\mathrm{mask}(c)$ after the children, apart from $c$'s own bit set at
its visit. So at the end bit $x$ of $\mathrm{mask}(c)$ is set iff $c = x$ or
bit $x$ is set in some child's final mask, which by the same induction is iff
$c = x$ or some child is in $\Rch{x}$, that is, iff $c \in \Rch{x}$. Each
commit is visited once and each parent edge relaxed once, with one
$\lceil m / w \rceil$-word OR per relaxation and one traversal step even when
$m = 0$, giving the cost.
\end{proof}

Children-before-parents is the reverse of the order \spec{25} requires of an
adapter, which returns parents first. Afterwards $c \in \mathrm{after}(b)$ is
one clear bit, a cancellation test is a non-empty intersection of
$\mathrm{mask}(c)$ with the cancel markers whose scope contains the package,
and the proper-ancestor test of \spec{13.4b} is one set bit together with $c
\ne x$. Under \spec{27} the pass runs once per repository over that
repository's own markers, with repository identity in every key
(\cref{def:exec-windows}). A pass confined to the union of the pending windows
and the parent edges inside it sets the same bits for the union's commits,
because a window is closed under descendants and no path leaves the union and
returns to it; what remains in $H + A$ is finding the union.

\begin{corollary}[the fresh-window identity]
\label{cor:fresh-identity}
With $\mathrm{after}(b) = \Rch{\Head} \setminus \Rch{b}$, a commit $c$ is in
$\Winf{\st}{p}$ exactly when $c \in \mathrm{after}(\stcommit{\st}{p})$ and $c
\in \mathrm{after}(\basec{\st}{p})$. Every window is therefore a function of
one boundary commit, a fresh membership test is the conjunction of two
lookups, and no reachability set is keyed by a $(\text{stable}, \text{fresh})$
pair; the identity holds whether or not the newest baseline descends from the
stable one.
\end{corollary}

\begin{proof}
By \cref{def:window}, $\Winf{\st}{p} = (\Rch{\Head} \setminus
\Rch{\stcommit{\st}{p}}) \setminus \Rch{\basec{\st}{p}} = \Rch{\Head} \cap
\overline{\Rch{\stcommit{\st}{p}}} \cap \overline{\Rch{\basec{\st}{p}}} =
\mathrm{after}(\stcommit{\st}{p}) \cap \mathrm{after}(\basec{\st}{p})$, with
an empty subtracted set where a boundary is undefined. Nothing in the
derivation relates the two boundaries to each other.
\end{proof}

Window storage may still be shared only inside one plan and one fixed
snapshot, between packages whose boundary identity is equal, and the shared
set must be immutable; a materialised fresh window is shared only when both
boundaries are equal. Storing $\Win{\st}{p}$ explicitly per package costs
$O(P \cdot C)$ memory; $k + 1$ dense bitsets of $H$ bits, or $O((k+1) \cdot
C)$ bits over the union once it is known, or a sparse $O(I_w)$, each make
admission a lookup and none turns the history traversal into $O(C)$.

\begin{proposition}[early stop, \spec{13.11}]
\label{prop:early-stop}
Let the pass of \cref{prop:marker-pass} visit commits in an order known to be
children-before-parents without a full walk, decreasing commit-graph
generation number for instance. Once every unvisited commit that has a visited
child carries all $k$ boundary bits, every unvisited commit lies in no
window and is named by no tuple of \spec{13.4}, so the pass may stop. The stop
is sound only when no package lacks a stable baseline, and only after every
commit named by an \code{Edits}, \code{Deletes} or \code{Reverts} footer has
been visited.
\end{proposition}

\begin{proof}
\Head{} has the greatest generation number among reachable commits, since the
number increases strictly along child edges, so \Head{} is visited first. Let
$c$ be unvisited and take any path $\Head = c_0 \to c_1 \to \dots \to c_j =
c$. It has a first unvisited commit $c_i$ with $i \ge 1$, whose child $c_{i-1}$
on the path is visited, so $c_i$ carries every boundary bit: $c_i \in \Rch{x}$
for every boundary $x$, and $c$, an ancestor of $c_i$, is in every $\Rch{x}$
too. Every train window is $\Rch{\Head} \setminus \Rch{x}$ for a stable
boundary $x$ when every package has one, every fresh window is a subset of a
train window, and every owed window of \spec{13.3} excludes $\Rch{x}$ for the
provider release $x$ that bounds it, which is among the $k$ boundaries; so $c$
is in no window. A package without a stable baseline keeps $\Rch{\Head}$ as
its window, which no boundary bit excludes, and the diagnostics \diag{E210}
and \diag{W213} of \spec{13.4b} read the ancestry of a correction target that
lies outside every window; each is the stated condition.
\end{proof}

\begin{remark}[the delivery test]
\label{rem:delivery-cost}
The delivery test of \cref{def:sources} is one window bit for a target that
has not released past the unit's commit, which is every target in the ordinary
case. A target that has costs one ancestry question per release of each source
past that commit, answered from the marker pass when those releases are
markers, and such targets exist only after a consumer got ahead of a provider.
The owed windows those targets add are counted in $k$, and are empty while no
consumer has got ahead.
\end{remark}

\subsection{Safe traversal reuse}

The superseded bucketing proposal grouped units by window class and directive
tuple and traversed from the union of their sources. That key is insufficient:
unit identity affects admission, cancellation, scope, inherited channel,
self-source exclusion and output provenance. \spec{13.11} identifies identical
graph-walk inputs as the one safe memoisation key, and additionally permits
one unbounded walk per (source package, edge kinds) whose union across the
sources of one unit, minus those sources, at the least level per target, is
exact within that unit. Other optimisations still require an equivalence
argument. The following historical result characterises why the naive merge
across units is unsafe.

\begin{definition}[the naive merged candidate set, \spec{13.11}]
For $S \subseteq \Pkgs$ let $\delta^{+}(S,d)$ be the least length of a chain of
\emph{at least one} edge from a member of $S$ to $d$, and $\reachset{S}{\ell}^{+}
= \{d : \delta^{+}(S,d) \leq \ell\}$. This is the naive positive-edge candidate set; \SPEC{} additionally requires attribution to an eligible unit and that unit's self-source exclusion.
\end{definition}

\begin{proposition}[limit of naive merged bucketing]
\label{prop:bucket}
Let $u_1,\dots,u_m$ share a depth $\ell$ and let $S_i = \srcp{\st}{u_i}$,
$S = \bigcup_i S_i$. Then
\[
  \bigcup_{i} \reachset{S_i}{\ell} \;\subseteq\; \reachset{S}{\ell}^{+},
\]
the two sides agree on every $d \notin S$, and they differ exactly on the
packages
\[
  X \;=\; \bigl\{\, d \in S \;:\; \delta^{+}(S,d) \leq \ell
      \ \text{ and }\ \forall i,\ d \in S_i \text{ or } \dist{S_i}{d} > \ell \,\bigr\},
\]
which $\reachset{S}{\ell}^{+}$ admits and no individual traversal does. In
particular, if no $\srcp{\st}{u_i}$ contains a package reachable from that same
set within $\ell$ edges, then $X = \emptyset$ and bucketing is exact.
\end{proposition}

\begin{proof}
Inclusion: if $d \in \reachset{S_i}{\ell}$ then $d \notin S_i$ and
$\dist{S_i}{d} \leq \ell$, so the witnessing chain has length at least $1$ and
starts in $S$, giving $\delta^{+}(S,d) \leq \ell$.

Agreement off $S$: let $d \notin S$ and $\delta^{+}(S,d) \leq \ell$. The
witnessing chain starts at some $p \in S_i$, and $d \notin S_i$, so
$d \in \reachset{S_i}{\ell}$.

The difference: for $d \in S$, membership of $\reachset{S}{\ell}^{+}$ says only
that some chain of length $1..\ell$ reaches $d$ from $S$, whereas membership of
$\bigcup_i \reachset{S_i}{\ell}$ requires some $i$ with $d \notin S_i$ and
$\dist{S_i}{d}\leq\ell$. Negating the latter gives the stated $X$.
\end{proof}

\begin{example}[the bucket that over-admits]
\label{ex:bucket}
On \cref{fig:workspace}, take the single unit
\[
  u \;=\; \code{fix(core,ui)\^{}major: x},
\]
so $S_1 = \{\code{core},\code{ui}\}$ and $\ell = 1$. Per \cref{ex:reach},
$\reachset{S_1}{1}$ excludes \code{ui}, which takes the unit's own \bpatch{}
directly. But $\delta^{+}(S_1,\code{ui}) = 1$ by the edge
$\code{ui} \to \code{core}$, so $\code{ui} \in \reachset{S_1}{1}^{+}$: the
merged rule gives \code{ui} the propagated \bmajor{} as well, publishing it two
levels higher than the specification says, with no diagnostic to show for it.
Here $X = \{\code{ui}\}$.
\end{example}

\begin{remark}[what $X$ costs, and what \spec{13.11} says about it]
\label{rem:bucket}
The warning box in \spec{13.11} is about the opposite error, subtracting the
merged source set $S$, which withholds the bumps of a package that is a source
of one unit and a legitimate target of another, and now requires a per-unit self-source test. Merely using $\reachset{S}{\ell}^{+}$ over-admits on $X$: a unit whose own scope-set
contains both a package and a dependent of it within the depth bound. Then the
dependent, which \spec{9.2} excludes from that unit's targets because the unit
bumps it directly, receives the unit's \emph{propagated} bump as well, a
difference whenever \code{Propagate} names a value above the unit's own bump,
as in \code{fix(core,cli)\^{}major}. \Cref{ob:bucket} preserves the 1.0.0 counterexample that motivated the current rule.
\end{remark}

\subsection{Hoisting}

\begin{proposition}[the resolvability predicate is loop-invariant]
\label{prop:hoist}
$\resolv{\st}{u}{d}$ as computed by the source scan
\[
  \exists p \in \srcp{\st}{u}:\ \chanres{\st}{p} = \stable \ \lor\
  \chanres{\st}{p} = \chanres{\st}{d}
\]
is equal, for every $d$, to the membership test of \cref{def:resolvable}
against the precomputed set $\srcch{\st}{u}$. Evaluating it once per unit and
testing membership per target costs $O(D + S)$ where the literal form costs
$O(D \cdot S)$; hoisting $\cof{u}$ and $\pscope{u}$ likewise takes the unit
from $O(D(S+\Sigma))$ to $O(D + S + \Sigma)$.
\end{proposition}

\begin{proof}
A set contains $\chanres{\st}{d}$ exactly when some $p \in \srcp{\st}{u}$ has
$\chanres{\st}{p} = \chanres{\st}{d}$, and contains \stable{} exactly when some
$p$ has $\chanres{\st}{p} = \stable$; the disjunction of the two membership
tests is therefore the disjunction of the two existentials, which is the source
scan. The set has one entry per \emph{distinct} channel among the sources.
The predicate reads $d$ only in that one comparison, and \cref{a:scope}
and $\cof{u}$ do not read $d$ at all, so all three may be computed before the
target loop.
\end{proof}

The cost matters because of where it lands: a unit carrying \code{\^{}\^{}}
across a wide scope-set admits thousands of targets from a source set of one or
two packages, and the literal reading re-walks those sources thousands of times
to re-derive an answer that never changed. Hoisting is not an optimisation to
be justified by profiling; it is what the predicate's own dependency structure
already says. The bound covers the hoisted scope and channel predicates only;
cancellation intersections cost up to $\mathrm{bw}(\mathrm{cancels})$ per
queried target, and materialising provenance charges the source insertions in
$Z_v$ even when deduplication leaves $Z$ unchanged.

\subsection{Inverted channel resolution}

\begin{proposition}[inversion]
\label{prop:invert}
Building the candidate lists by pushing from units to the packages their
scope-sets resolve to yields, for each $p$, the same list as scanning the
window for directives naming $p$, provided units are pushed in \spec{11.6}
precedence order. It costs $O(U + I + P)$ against the literal $O(P \cdot U)$; the incidence term cannot be omitted because those appends are output work.
\end{proposition}

\begin{proof}
A package $p$ appears in the pushed list of $u$ exactly when $u$ carries a
\code{Channel} directive, $p \in \scopef{u}$, and $\cof{u} \in \Winf{\st}{p}$,
which is the pulled form's membership condition read in the other
direction, so the two lists have the same elements. They have the same order
because the push visits units in the order the pull would sort them into.
The push performs one pass over units, each contributing $|\scopef{u}|$
appends against scope resolutions \spec{13.4} has already paid for, and one
flat pass over packages.
\end{proof}

Two obligations come with the inversion and neither is optional. Units must be
pushed in \spec{11.6} order, or the first-match-wins loop below silently picks
a different directive, a determinism bug that surfaces only when one package
is named by two commits. And the \diag{W186} count is over the whole candidate
list, not the prefix examined before the winner was found, so the list must be
complete before it is read; an implementation that fuses the two passes and
evaluates a package as soon as its first candidate arrives loses \diag{W186}.

\subsection{Deterministic order and failure closure}

\begin{proposition}[the publish heap]
\label{prop:heap}
The order of \cref{def:order} is computable in $O(E + P \log P)$ with a
comparison heap over the ready set, against $O(P^2 + E)$ for scanning the
ready set for its least name at every step. No comparison-based structure
improves the bound for unsorted names: a workspace with $P$ unsorted names and
no edges makes the sequence a sort of its names.
\end{proposition}

\begin{proof}
Kahn's algorithm~\cite{kahn} visits each package and relaxes each edge once;
with the zero-indegree packages kept in a heap keyed by name, each package
enters and leaves the heap once at $O(\log P)$ and each relaxation is $O(1)$,
giving $O(E + P \log P)$. A linear scan of the ready set at each of $P$ steps
costs $O(P)$ each. For the lower bound, with no edges every package is ready
at once and \cref{def:order} returns the names in byte-wise order, so any
method that decides the order by comparisons performs a comparison sort of
$P$ keys, which needs $\Omega(P \log P)$ comparisons. This is a worst-case
statement for unsorted comparison keys and not an instance-optimality result:
where name ranks or sorted names are already provided, the reduction proves
no new lower bound, and any other structure must provably return the same
least element at every step.
\end{proof}

The historical bound of $O(P + E)$ for an arbitrary topological order predates
the byte-wise tie-break \spec{19.2} now requires; \cref{prop:order} proves
that the required order is well defined and machine independent, not that it
is linear.

\begin{proposition}[the blocking closure]
\label{prop:closure}
The blocked set of \cref{def:blocking} is computable in $O(P + E)$ traversal
work per run, plus one read of each planned package's recorded causes, by
adding each failure to an incremental multi-source reverse traversal over
\code{publish.blockingKinds} on the full workspace graph and marking each
newly reached consumer once. The literal \code{run} of \spec{19.3}, which
tests each planned package by a fresh transitive walk, is $O(P \cdot (P +
E))$ whatever the number of failures, and a walk from each failure is $O(F
\cdot (P + E))$.
\end{proposition}

\begin{proof}
A package transitively depends on a failed or blocked package exactly when it
transitively depends on a failed one, since a blocked package itself depends
on a failed one; so the set \code{bad} of \code{run} is non-empty for $d$
exactly when $d$ is reached from $F$ over the reversed edges. Marking a
package on first reach and never expanding it twice visits each package and
each reversed edge at most once across all failures, which is $O(P + E)$
however the failures interleave with publication. Whether a reached package is
blocked or proceeds is then the own-cause test, which reads the causes the
plan already retains for that package and excludes those from failed or
blocked providers; packages are examined in publish order, so every
provider's status is known when its dependent is examined. Recording every
failed ancestor or path adds output-sized work, and the first cause and the
\diag{W194} order still follow the deterministic publish and failure order
rather than queue or hash iteration order.
\end{proof}

\begin{proposition}[build readiness in linear time]
\label{prop:readiness-cost}
The build-stage edges of \cref{def:taskgraph} are computable in $O(P + E)$:
give every package that does not build a pass-through node, direct one relay
edge from provider to consumer for every dependency edge not declared
\code{none}, and take the build constraints as reachability in that relay
graph. A search per building pair through the packages between them is $O(P
\cdot (P + E))$. After transitive reduction at most two task edges per
dependency edge survive.
\end{proposition}

\begin{proof}
A path in the relay graph from $P$'s node to $C$'s node is exactly a
dependency path from $C$ to $P$ with no \code{none} edge, since relay edges
exist for precisely the edges not declared \code{none} and pass-through
nodes relay without constraint; so $\exbuild{P}$ precedes $\exbuild{C}$ in the
relay graph exactly as \cref{def:readiness} requires. Building the relay graph
touches each package and each edge once. Scheduling then follows the relay
edges, which the executor already does for the publish stage, so no
per-pair search is needed. Alternatively a memoised index of the nearest
building packages each package reaches without a \code{none} edge is built in
one visit per package in reverse topological order, at the cost of the set
unions it stores, which are output-sized. For the last claim, each dependency
edge $C \to P$ contributes at most one build-stage constraint (build-to-build
or publish-to-build, by its relation) and one publish-stage constraint, and a
constraint carried through intermediate packages is implied by the relayed
edges, so the reduction keeps at most two per dependency edge.
\end{proof}

\subsection{The profile rows}

The polyrepository rows are sums over repositories and deliberately not
globally linear in fleet history: distinct boundaries can require distinct
reachability sets, output can contain $P \cdot U$ scope incidences, and
propagation can still materialise $Z$ contributions. Snapshots: the control
gitlink transitions are indexed in one pass per fixed control snapshot and
each source snapshot is walked once, $O(G + \sum_q (H_q + A_q))$, never once
per consumer. Windows: one marker pass per repository, $O(\sum_q (H_q +
A_q)\,\mathrm{bw}(m_q) + I_w)$, with the boundaries of \cref{def:exec-windows}
counted per repository in $K_q$ as revisions and not as $(\text{stable},
\text{fresh})$ pairs. Publication input closure: the graph augmented with
shared-version-group membership can contain cycles, so its strongly connected
components are condensed in $O(P + E + V)$ and visited in reverse topological
order with one bitset union per condensed edge, $O((P + E + V)\lceil Q /
w\rceil)$; a dense representation costs $O(P \lceil Q / w\rceil)$ words and an
expansion into name lists up to $O(P \cdot Q)$. Link evidence and settlement:
the link tree is rooted once in $O(Q)$, each (repository, revision) tree is
read at most once per plan, and a package settles with one commit per
recording repository, at most $Q - 1$ and two in a star. Record comparison:
one inventory per repository under the lock, $O(T)$ each
(\cref{def:records-lock}). Every index and cache key preserves repository
identity, no cache synthesises ancestry between repositories, and ancestry
and walk caches are bounded by configured memory or an eviction policy over
all their tiers, because a cache of every queried pair grows quadratically in
$H_q$; recomputation after eviction preserves semantics and may lose the
cached time bound, a trade-off an implementation states.

The execution rows count tasks and nodes. Task graph and readiness: a
counter of unfinished predecessors per task, decremented once per completed
predecessor, makes readiness $O(T_k + D_k)$ over the run, where rescanning
every task after each completion is $O(T_k \cdot (T_k + D_k))$. Placement:
choosing a compatible free node per task is $O(T_k \cdot W_n)$, and $W_n$ is
small enough that a free list per compatibility class is rarely worth its
bookkeeping. Output transfer: one transfer per distinct (output set, node)
pair bounds the count by $P_o \cdot W_n$ where a transfer per consumer task
would reach $T_k \cdot P_o$, and verified bytes are not fetched again
(\cref{def:exec-outputs}). Coordination polling: the branch-name prefix of
\cref{def:exec-branches} lets a node list only the refs addressed to it and
fetch once per changed tip, where listing every live ref is $O(L_r)$ per tick,
and the refs of dead runs tax every tick until removed. None of these rows
bounds transfer bytes or command durations, which \cref{sec:execution} treats
separately.

\subsection{Space}

Memory is counted in machine words, except input and output byte buffers,
with indexed inputs retained.

\begin{center}
\footnotesize
\begin{longtable}{@{}L{4.2cm}L{4.0cm}L{6.4cm}@{}}
\toprule
Representation & Retained or auxiliary space & Qualification \\
\midrule
\endhead
Package and history adjacency & $O(P + E + H + A)$ & input indexes; not repeated per package \\
Parsed messages and resolved incidences & $O(N + U + I)$ & byte buffers plus unit and incidence records; cached resolution has its own key \\
Dense history marker rows & $O(H \lceil m / w \rceil)$ & plus traversal state; zero markers need no rows \\
Explicit package windows & $O(I_w)$ & optional materialisation; shared boundary indexes can avoid it \\
Cancel masks & $O(P_c \lceil \mathrm{cancels} / w \rceil)$ & plus scoped-package lookup and diagnostic state \\
One propagation BFS & $O(P)$ auxiliary & on top of shared adjacency; retained walk caches have a separate budget \\
Direct channel candidates & $O(P + I)$ as lists, or $O(P)$ summaries & summaries require the deferred-output rule of \spec{13.8}; diagnostics and provenance are additional \\
Contribution and provenance records & $O(Z)$ & streaming reduces retention only where later phases and output permit \\
Heap publish order, blocking membership & $O(P)$ auxiliary & shared graph excluded; all-cause or path diagnostics add output-sized storage \\
Publication input closure bitsets & $O(P \lceil Q / w \rceil)$ & shared component sets may reduce storage; expanded lists cost up to $O(P \cdot Q)$ \\
Link evidence cache & retained pins plus keys & $B$ counts trees, not their entries or bytes; dense hub revisions can be large \\
\bottomrule
\end{longtable}
\end{center}

No implementation should add every alternative in this ledger by default.
Peak simultaneously live allocations, construction scratch space, immutable
input buffers, retained output and cache keys are what count; batching the
marker computation limits scratch space but does not bound retained marker
columns, and a transposed per-marker layout has the same $H \cdot m$ logical
bits as the row layout with neither universally better. Linear scanning is
worst-case time-optimal when every input byte must be inspected, explicit
output needs time at least its size, a single adjacency traversal and the
blocking closure have tight linear bounds, and the publish heap matches the
comparison-sorting lower bound for unsorted names. These are conditional,
phase-specific statements: neither $\Omega(Z)$ output work nor a dense $H
\times m$ or $P \times Q$ representation proves that every instance requires
that representation or that the complete planner is jointly time- and
space-optimal.

\subsection{Parsing}

Lexical parsing and resolution have distinct cache keys. Lexing is pure in the
repository identity, immutable record identity, parser and schema identity and
parsing configuration. Scope and changed-path resolution additionally depends
on the workspace package and path index, graph and resolution configuration. A
lexical cache hit therefore does not cache $R$ or $I$, and record immutability
does not imply a whole-plan cache. File ownership is a longest-prefix probe of
the path's own ancestor directories, a literal scope term is one map lookup, a
glob whose only \code{*} is its last byte selects a contiguous range of the
sorted names in $O(\log P)$ plus its matches, and a scope-set depends on its
commit only through \code{.} and, under \spec{27}, through the owning
repository, so it may be resolved once per distinct text; none of this changes
$I$, which is output. Vectors 82h to 82j exercise direct-channel resolution,
inherited-channel provenance and per-unit admission; they do not exhaustively
validate these cache keys.

\begin{proposition}[linear parsing, \spec{20.7}]
\label{prop:parse}
A conforming parser reads a message of $n$ bytes in $O(n)$ time and $O(1)$
auxiliary space beyond the parsed result, with no input causing superlinear
behaviour.
\end{proposition}

\begin{proof}[Proof sketch]
\spec{20} specifies every step as a forward index scan over the byte array with
no backtracking: separator detection compares a fixed token at line starts;
header parsing advances a cursor through type, scope-set, inline directives,
\code{!} and \code{:~}, each transition determined by the current byte and a
bounded amount of already-scanned state; footer detection tests each line
against a fixed set of key shapes. Each byte is examined a bounded number of
times, and the working state is a fixed number of indices and flags. No
construct requires re-examining a prefix, which is the property that
distinguishes it from a backtracking regular-expression engine and is why
\spec{18.3} can promise bounded work on untrusted commit messages.
\end{proof}

This is the one bound in the table stated for the parser rather than the
planner, and it is here because \spec{18.1}'s threat model includes commit
messages authored by anyone who can open a pull request. A superlinear parser
would make that a denial-of-service surface in CI; \cref{prop:parse} is why the
specification prescribes an index scan and relegates regular expressions to an
appendix.

\subsection{What none of this may change}

Every representation above is internal. The plan, the diagnostics and their
order must be identical to the literal reading (\spec{17.2},
\cref{thm:determinism}), and an implementation that trades a different plan
for speed does not conform however fast it is. Appendix B.7a of \SPEC{} tests
exactly the transformations that are almost equivalence-preserving, a hoist
that collapses a multi-channel source set, an inversion that loses
\spec{11.6} order, a fused pass that miscounts \diag{W186}, a skipped channel
pass that leaves a channel unset; it is not exhaustive, and an optimisation it
does not name is still the implementer's to justify.
